
Paraphrasing defeats contamination detection. When Llama-2-13B is trained on rephrased MMLU items, it achieves 85.9\% accuracy while evading 13-gram overlap filters \citep{yang2023rethinking}. This exposes a fundamental limitation: surface-form matching cannot detect contamination that preserves semantics but varies phrasing.

Current detection methods occupy two extremes. N-gram overlap, established by GPT-3's 13-gram threshold \citep{brown2020language} and refined by GPT-4's 50-character threshold \citep{openai2023gpt4}, scales to trillion-token corpora but fails completely on paraphrased contamination. Quiz-based methods such as DCQ \citep{golchin2023data} achieve higher memorization detection by probing model behavior, but provide only binary outputs and require per-item evaluation, limiting scalability. Neither approach offers a continuous-valued, scalable, paraphrase-resistant contamination metric.

This work hypothesizes that contamination produces representation-level changes detectable via confidence behavior, regardless of surface form. The intuition follows from learning theory: a model trained on benchmark items in various phrasings should develop robust semantic representations that generalize uniformly across paraphrase variations. This uniformity should manifest as consistent confidence scores---low variance---across paraphrases of contaminated items. The Semantic Saturation Index (SSI = 1/variance) is proposed, where high SSI indicates potential contamination.

The investigation validates the underlying mechanism but exposes the metric's failure. Through controlled experiments on Mistral-7B with known contamination levels (0\%, 10\%, 50\% of MMLU), the following is demonstrated:

\begin{enumerate}
\item Contamination injection creates measurable training exposure, with 31.1\% accuracy gain at 50\% contamination (H-M1).
\item Paraphrase-augmented training produces representation invariance, with Mean Pairwise Similarity (MPS) differing by 0.065 between training conditions (Cohen's d=0.52, p=0.008) (H-M2).
\item Representation invariance correlates with confidence uniformity (Pearson r=$-0.517$, Cohen's d=0.58) (H-M3).
\end{enumerate}

However, the SSI metric itself fails to discriminate contamination on real inference data. On actual MMLU inference with Mistral-7B, SSI achieves AUC=0.506---essentially chance-level classification (H-M4). The 1/variance formulation amplifies noise: standard deviations exceed means at all contamination levels, and SSI distributions overlap completely.

This paper makes three contributions:

\begin{enumerate}
\item The first empirical demonstration of the causal chain linking contamination to representation invariance to confidence uniformity.
\item Documentation of the failure of inverse-variance metrics for contamination detection, guiding future work toward alternative formulations.
\item Identification of a simulation-reality gap where mechanism hypotheses passed on simulated data but real-data evaluation revealed metric failure.
\end{enumerate}
